problem:  
Let \( (S,g_E) \) be a simply connected **standard Einstein solvmanifold** with solvable Lie algebra \( \mathfrak{s} = \mathfrak{a} \oplus \mathfrak{n} \), Einstein derivation \( D_E:\mathfrak{n}\to\mathfrak{n} \), and Einstein constant \(\lambda_E<0\).  Denote by \( \mathfrak{m} = \mathrm{Sym}^2(\mathfrak{s}^*) \) the space of left‑invariant symmetric bilinear forms on \(\mathfrak{s}\).  

Let \( E_{\mathrm{alg}}(g) = \mathrm{Ric}_{\mathrm{alg}}(g) - \lambda_E g \), with \( \mathrm{Ric}_{\mathrm{alg}}(g) \) computed in terms of the structure constants of \(\mathfrak{s}\), and let \( L_{\mathrm{alg}} = D E_{\mathrm{alg}}(g_E): \mathfrak{m}\to\mathfrak{m} \) be its linearization.  Trivial directions arise from scaling and derivations:
\[
T = \operatorname{span}\{\, h_0:=g_E,\, h_D(X,Y)=g_E(DX,Y)+g_E(X,DY): D\in \mathrm{Der}(\mathfrak{s}) \,\}.
\]

Explicit curvature calculations for Einstein solvmanifolds show \(L_{\mathrm{alg}}\) is diagonalizable when expressed relative to the eigenspaces of \(D_E\).  Define the **reduced kernel**
\[
K_{\mathrm{red}} = \{\, h\in \ker L_{\mathrm{alg}} : h \perp T \,\}.
\]
While in all examples known to date \(K_{\mathrm{red}}=\{0\}\), a general proof of this vanishing is not known.

> **Problem (Reduced‑Kernel Rigidity Conjecture for Einstein Solvmanifolds).**  
> Does every standard Einstein solvmanifold satisfy \( K_{\mathrm{red}} = \{0\} \)?  
> Equivalently, are all infinitesimal algebraic Einstein deformations of a solvmanifold induced solely by derivations and scaling?  

In the event that a nontrivial \(h\in K_{\mathrm{red}}\) exists, the quadratic term \(Q_{\mathrm{alg}}(h,h)\) in the expansion  
\[
E_{\mathrm{alg}}(g_E+h)=L_{\mathrm{alg}}(h)+Q_{\mathrm{alg}}(h,h)+O(\|h\|^3)
\]
is automatically nonzero (by known curvature identities) but might still be annihilated modulo \(T\). Thus a counterexample would not only provide a non‑rigid Einstein solvmanifold but would imply the existence of a **genuinely new algebraic Einstein deformation mechanism.**

Why is it a "good" problem:  
1. **Precision and conceptual clarity:**  
   The statement \(K_{\mathrm{red}}=\{0\}\) is perfectly sharp and verifiable: either the kernel of the linearized algebraic Einstein operator equals the subspace generated by derivations and scaling, or it does not.  It turns the broad rigidity expectation into a single concrete algebraic condition on the solvable Lie algebra, isolating the genuine undeformed content of Einstein solvmanifolds.

2. **Profound significance:**  
   A positive answer would establish universal *algebraic infinitesimal rigidity* for Einstein solvmanifolds, resolving a long‑standing implicit assumption underlying Heber’s uniqueness theorem.  
   A negative answer would produce—for the first time—a new family of Einstein solvmanifolds with nontrivial deformation directions. Either outcome would reshape the structural understanding of noncompact homogeneous Einstein spaces.

3. **Bridge between algebra and geometry:**  
   The problem explicitly relates the algebraic operator \(L_{\mathrm{alg}}\), derived from Einstein curvature equations, with the Lie‑theoretic object \(\mathrm{Der}(\mathfrak{s})\).  It unifies methods from differential geometry (Einstein equation linearization), Lie‑algebra theory (derivations, eigenvalue structures), and spectral algebra (diagonalizability and kernel analysis).  

4. **Extreme simplicity with deep reach:**  
   The question is short, self‑contained, and purely algebraic:  
   *“Does \( \ker L_{\mathrm{alg}} \) reduce to derivation‑ and scaling‑induced directions for every Einstein solvmanifold?”*  
   Yet resolving it would have major implications—from classifying Einstein solvmanifolds’ moduli to testing the deeper algebraic structure implicit in Lauret’s correspondence between Einstein metrics and nilsolitons.

Thus, this problem achieves *mathematical exactness, conceptual centrality*, and *broad potential impact*—qualities that define an outstanding, research‑level “good” problem.